\chapter{Implementation}\label{chap:chap4}
% Tenho de prv dar um reframe a isto pq "by building a visualization example from the ground up" é o que está a acontecer...
In this chapter, we will explore the inner workings of MIRVisualScore through the lens of practical code examples and patterns drawn from this project's example scripts. We provide simplified code snippets to show architectural patterns, and workflows that demonstrate how the codebase is structured, how its components function, and how it is intended to be used in practice. We will progress through the essential topics: setting up the required dependencies, understanding and obtaining the required input files, examining how the layer-based visualization system works, and finally, exploring how to compose multiple data layers into aligned score visualizations. Through this approach, readers should gain a clear understanding of the project's core components and workflows, enabling them to adapt and extend the codebase for their own applications.

\section{Dependencies and Setup}
% DEMASIADO LOW LEVEL
The depandancies and requirements and how to install them are laid out in the README.md of the codebase. Yet, we will go a bit more in depth in this section. MIRVisualScore as of now, runs using python 3.12.3, using other updated python environments might work, but this was the version used during the writing of this thesis. Using other python versions might cause conflicts with the rest of the codebase. For simplicity, we will run the codebase as if we were running it on a Linux machine, for running other operating systems, simply follow along as the README.md specifies. 

% DEMASIADO LOW LEVEL
\noindent We advise running this codebase with a virtual environment to avoid affecting other projects and repositories you might be working on your system.

% DEMASIADO LOW LEVEL
\noindent Other then the python version the rest of the codebase requires the following dependancies:
\begin{itemize}
    \item modusa==6.3.7
    \item matplotlib==3.10.9
    \item numpy==2.4.4
    \item librosa==0.11.0
    \item soundfile==0.13.1
    \item torch==2.11.0
    \item beat\_this==1.1.0
\end{itemize}

To set this up, as the README.md explains, first make sure you have python 3.12.3 installed in your system. Then you can simply run the following commands in the terminal:

\begin{verbatim}
# Create virtual environment
python3.12 -m venv venv

# Activate the virtual enviroment
source venv/bin/activate

#Install requirements
pip install --upgrade pip
pip install -r requirements.txt

\end{verbatim}


\section{Required Files}
For generating visualizations with audio and data aligned with the score we will, at least, three files: A audio file in .wav format with the performance recording, the warped score generated by ScoreWarp, the MAPS file, the .mei score and an .npz with BT data.

\subsection{Getting the Warped Score}
The warped score can be generated using the \href{https://iwk-digital.github.io/scorewarp/}{ScoreWarp online demo}.  In it, one needs to upload the score in .mei format and the MAPS file in order to download the score in svg format. 

The .mei file that contains the score in xml format. The easiest way to get a .mei score is to get the score in .musicxml format and then convert it to .mei, this can be done in multiple ways we will describe two: The first being through the \href{https://editor.verovio.org/}{verovio online editor}. The second, by using musescore's software that allows to then convert to .mei. In either way, it is important to get the original score in .musicxml format to be able to then convert it to .mei using the tools we described. 

\subsection{Getting the MAPS file}
The MAPS file, describes the onsets of the performance and links them to each corresponding score elements of the .mei score through their @xml.id . Getting the maps file might be a bit more cumbersome, but we can describe two different ways of obtaining it: 

The first is the one described in \citet{goeblLetsScoreWarpAgain2025a}, using the MIDI-to-MIDI matching algorithm by \citet{nakamuraPERFORMANCEERRORDETECTION}. However,  this method requires , along side with the .mei score, the performance in midi format. Using examples from the asap dataset \citep{peterAutomaticNoteLevelScoretoPerformance2023} will facilitate this process but if a midi performance is not possible to obtain then the maps file needs to be generated independently.

The MAPS file is at it's core a .json file, and for the purposes of the current MIRVisualScore version, it needs to provide the score's and performance onset data in the following manner:

\begin{verbatim}
{
    "obs_mean_onset": 0.853333333,
    "xml_id": ["x1iw1eq1"],
    "obs_num": 1
}
\end{verbatim}

\noindent Where: 

\begin{itemize}
    \item ``\textit{obs\_mean\_onset}'' is the time in seconds of a given onset.
    \item ``\textit{xml\_id}'' is the @xml.id of the corresponding onset note in the .mei score. In the case that the onset has more than one note, like a chord hit, this attribute can have more than one @xml.id. 
    \item ``\textit{obs\_num}'' acts as a simple counter for each detected onset. 
\end{itemize} 

\noindent MIRVisualScore relies on MAPS in two complementary ways: directly for onset visualization, and indirectly through ScoreWarp, which requires MAPS to generate the warped score aligned with the timeline SVG element which constitutes the alignment reference for all overlaid data.

\noindent ScoreWarp performs best when MAPS contains dense onset annotations across the full performance. However, it can still align scores when annotations are sparse by interpolating between available onset and their @xml.id anchor points. In practice, this allows different precision levels (i.e., downbeats, phrase boundaries, or measure-level anchors), trading annotation effort for alignment precision. Even minimal anchors (first and last onset) may be usable for short performances or segments, but intermediate alignment becomes less reliable due to interpolation error.

% =========== AI Recomendation ===========
% This provenance paragraph is probably too implementation-specific for the thesis; keep the fact that MAPS can be derived from alignment tools, but trim the workflow details about PP, exported txt files, and the helper function.
% ↓↓↓↓↓↓↓↓↓↓↓↓↓↓↓↓↓↓↓↓↓↓↓↓↓↓↓↓↓↓↓↓↓↓↓↓↓↓↓↓↓↓↓↓
\noindent For the examples that we provide in the making of our codebase we modified PP to obtain this data. PP's ``Piano Aligner'' plugin \citep{jiangPIANOPRECISIONADVANCING2025} matches performance onsets to score elements defined by the .mei file. In PP these are displayed in a regular SV time instance layer. We made a simple modification where, this layer, instead of displaying each onset with it's corresponding score position, ``\textit{measure + position\_in\_measure}'', it would display the @xml.id of each onset. With this we could then export the time instances as .txt files with ``\textit{obs\_mean\_onset}+\textit{xml\_id}'', then convert it into the MAPS format using ``\textit{TXT\_to\_Maps()}'' function in @src/functions/warp\_score.

\subsection{Providing BT data}
Since contemporary BT algorithms are ML based, they primarily function by generating a beat activation function, which is a continuous curve representing the probability of a beat occurring at any given frame. Visualizing this underlying probability curve, rather than just the discrete beat assessments, allows us to see the algorithm's internal confidence level, giving crucial insights on the inner workings of the model.

% =========== AI Recomendation ===========
% This is still a bit code-adjacent because it introduces the exact helper function and NPZ field contract; keep the concept of beat-confidence visualization, but condense the file-format explanation to the minimum needed for the thesis.
% ↓↓↓↓↓↓↓↓↓↓↓↓↓↓↓↓↓↓↓↓↓↓↓↓↓↓↓↓↓↓↓↓↓↓↓↓↓↓↓↓↓↓↓↓
\noindent In @src/functions/Beat\_Layers, the function ``\textit{Run\_BeatThis()}'' generates an .npz file that contains the required data for this to be later visualized. This particular function creates the .npz file generated through BeatThis. However, the .npz file format can serve as an intermediate for BT visual layers, as long as it is structered like it is described in ``\textit{Run\_BeatThis()}''. The .npz should contain five key components from the BT model:

\begin{table}[h]
\centering
\begin{tabular}{ll}
\hline
\textbf{Field} & \textbf{Purpose} \\
\hline
\texttt{beat\_times} & Time coordinates (seconds) aligned with probability frames \\
\texttt{beat\_probs} & Continuous beat activation function values \\
\texttt{downbeat\_probs} & Continuous downbeat activation function values \\
\texttt{detected\_beats} & Discrete beat positions derived from peak detection \\
\texttt{detected\_downbeats} & Discrete downbeat positions from peak detection \\
\hline
\end{tabular}
\end{table}

\noindent The probabilistic fields capture the algorithm's internal confidence across time, while the detected positions represent the final beat assessments. This separation enables visualization of both the continuous confidence landscape and discrete beat estimates.

\subsection{Summarizing Required Files}
We now know the purpose of each required file and how to get it. To summarize all the files that are directly involved in the example, we are constructing:

\begin{table}[h]
\centering
\begin{tabular}{ll}
\hline
\textbf{File} & \textbf{Description} \\
\hline
\texttt{.wav} &  File with the audio recording of the performance \\
\texttt{.mei} &  XML Score of the musical piece\\
\texttt{.maps.json} &  Contains data that links onsets with score elements of the .mei\\
\texttt{.svg} &  The score in visual format generated by ScoreWarp (requires .mei and .maps \textit{a priori})\\
\texttt{.npz} &  File with algorithmic BT data\\
\hline
\end{tabular}
\end{table}

\section{Calling Modules, Methods and Functions}
The codebase divides it's functionalities into 4 main categories, where each, has it's own dedicated python script. They are located inside @src/functions and are:
\begin{itemize}
    \item \texttt{\textbf{visualization\_system.py}}: Focused on establishing the visualization methods for rendering the visual elements, and for managing svg outputs. 
    \item \texttt{\textbf{Spectrogram\_layer.py}}: Responsible for rendering spectrogram visualizations of audio.
    \item \texttt{\textbf{Beat\_Layers.py}}: Dedicated for classes and methods that are involved in the visualization of beat data.
    \item \texttt{\textbf{warp\_score.py}}: For methods that play a rule in the score-to-audio alignment between the warp svg score and the other elements that can be displayed along side the score (the ones generated by the other categories). 
\end{itemize}

\noindent In \texttt{\textbf{visualization\_system.py}}, the  ``\texttt{class Layer(ABC)}'' defines the abstract base class for all visual components. This class provides a common interface and structure for creating various visualizations, ensuring consistency and reusability across the codebase. The Layer class defines methods that layer methods should implement: \texttt{load\_data()}, \texttt{draw()}, and \texttt{to\_svg\_group()}. 

% =========== AI Recomendation ===========
% This paragraph is useful, but it is close to source-code documentation; compress it to a short module-level role statement and avoid naming every method unless one of them is central to the thesis argument.
% ↓↓↓↓↓↓↓↓↓↓↓↓↓↓↓↓↓↓↓↓↓↓↓↓↓↓↓↓↓↓↓↓↓↓↓↓↓↓↓↓↓↓↓↓
\noindent Each of these three methods serves a distinct purpose in the visualization pipeline:

\begin{itemize}
    \item \texttt{load\_data(**kwargs)} — Loads and validates all necessary data for the layer. For instance: \texttt{class Spectrogram(Layer)} computes the mel-scale spectrogram from an audio file, \texttt{class BeatProbabilityLayer(BeatLayer)} loads BT data from the .npz file, and \texttt{Onsets\_Layer} extracts onset times from a MAPS JSON file.
    \item \texttt{draw(ax, shared\_data)} — Renders the visualization onto a matplotlib axis object. It accesses a \texttt{shared\_data} dictionary that contains context information (such as audio duration, file paths, and SVG coordinate mappings) these are shared across all layers. This method returns a tuple of (lines, labels) which are collected for the legend.
    \item \texttt{to\_svg\_group(shared\_data)} — Converts the rendered visualization into vector-based SVG markup. This is essential for creating composited visualizations that can be embedded within the score SVG. Not all layers support SVG export; those that do, return SVG group markup, while others return \texttt{None}.
\end{itemize}

\noindent To orchestrate multiple layers into a unified visualization, the \texttt{Visualizer} class acts as a central coordinator. It maintains a collection of \texttt{Layer} objects and manages the workflow:
%Talvez possa ser demasiado tmb
\begin{itemize}
\item \texttt{add\_layer(layer)} — Appends a layer to the visualization
\item \texttt{load\_all\_layers(**kwargs)} — Sequentially loads data for all registered layers, and when an audio path is provided, extracts the audio duration and stores it in \texttt{shared\_data}
\item \texttt{draw()} — Renders all layers onto a single matplotlib figure, collects lines and labels for the legend, and enforces consistent time-axis limits across layers
\end{itemize}

\noindent The \texttt{shared\_data} dictionary serves as a coordination mechanism between layers and export steps. It is initialized empty and populated during execution with values such as the audio duration, file paths, axis objects, and SVG conversion context. This lets different parts of the visualization pipeline to work together by using shared data, without needing to depend on each other directly.

\noindent The composition pattern is best illustrated through a practical example. The following code creates a visualization combining differnt audio data layers:

% =========== AI Recomendation ===========
% Keep the example, but if the thesis version becomes long, trim the code walk-through to the minimum needed to show the pipeline rather than describing every parameter or line.
% ↓↓↓↓↓↓↓↓↓↓↓↓↓↓↓↓↓↓↓↓↓↓↓↓↓↓↓↓↓↓↓↓↓↓↓↓↓↓↓↓↓↓↓↓
\begin{verbatim}
#Initialize a visualizer instance. 
example1=Visualizer()   

# Add layers to the visulization. 
example1.add_layer(Onsets_Layer(onset_color=(1, 1, 1), 
                                line_Width=0.5))
example1.add_layer(BeatAccurateLayer(beat_color=(1, 0, 0), 
                                downbeat_color=(0, 0, 1)))
example1.add_layer(BeatProbabilityLayer(color=(1, 0, 0)))

#Load data for layers.
example1.load_all_layers(audio_path  =audio_file_example1,
                         maps_file   =maps_file_example1,
                         beat_file   =beat_example1)
                            
#Render the visualization:
fig, ax = example1.draw()

#Show visualization in Matplotlib window:
plt.show()
\end{verbatim}

\noindent The above example illustrates the layer-based implementation philosophy. At its core, the codebase composes matplotlib visualizations and arranges them so they can be displayed alongside a musical score and audio representations. The project relies on the SVG format because Verovio (and subsequently ScoreWarp) render the score in SVG, besides, embedding external SVG score markup into matplotlib plots is not straight forward especially given the complexity of musical scores. Because SVG is vector-based, 2D graphics can be exported without losing resolution. However, some data types, (e.g. spectrograms), cant be presented in this way, thus we employed a way to render plots as PNG's as well. This trade-off motivated the focus on SVG format. This preserves score elements and most data layers with high fidelity, even tough certain elements might need to be rasterized.\\

% =========== AI Recomendation ===========
% This is the right level of argument, but the explanation becomes repetitive if you keep expanding every export step; keep the SVG-versus-PNG trade-off and cut the rest into one compact pipeline summary.
% ↓↓↓↓↓↓↓↓↓↓↓↓↓↓↓↓↓↓↓↓↓↓↓↓↓↓↓↓↓↓↓↓↓↓↓↓↓↓↓↓↓↓↓↓
\noindent Matplotlib generated layers have to afterwards be rendered into SVG. In the case that the layers do not require to be rasterized, by calling \texttt{Visualizer().TurnLayersIntoSVG()} we create an SVG render of the plots. In doing this, each layer is made into its own self-contained SVG <g> element in the xml structure, this is done automatically through the \texttt{to\_svg\_group()} method.\\

\noindent In the case of a layer that needs to be rasterized, i.e., a spectrogram, the plot is rendered into a PNG using \texttt{Visualizer().TurnLayersIntoPNG()}, so that then, it can be embedded into a SVG using \texttt{Visualizer().Align\_Score\_and\_PNG()}. This method takes care of embedding the PNG into the SVG and aligning it with the score.\\

\noindent Afterwards a visualization in SVG can combine the rendered elements into a single combined SVG, for this we call \texttt{Visualizer().combine\_AlignedScore\_with\_Layers()} which takes care of merging the score and all imported layers into a single SVG file. In it, the layers (including the rasterized ones) are embedded into the same xml class Layers keeping the XML structure previously generated through \texttt{Visualizer().TurnLayersIntoSVG()}.\\ 

\noindent Finally, a complete multi-visualization composite can be rendered through \\ \texttt{Visualizer().create\_final\_SVG()}. This method combines different provided SVG files into a single one. It keeps the same XML structure philosophy, this time, one level above, i.e., each previously rendered SVG keeps the same structure and are treated as independent children of a new added root. In this final step, visualization can be offset vertically to allow for side-by-side comparisons.\\

\noindent In a sense, the codebase can be characterized by this layer-based philosophy, where layers are visual elements that can be combined and arranged. The project does this at two different levels, the first is at its core a matplotlib plot generator, and the second treats SVG elements as upper-level visual layers where they can be combined and arranged as individual elements. 

\section{Creating the Audio Visualizations}
% Core narrative: 
%   Show how each layer type (Spectrogram, Beat, Onset) works internally and how they integrate into the Visualizer.
% Key Topics:
%   Layer-based visualization philosophy -- explain that the audio visualizations are built as reusable layers rather than one-off plots.
%   Spectrogram as a time-frequency view -- describe its role as the audio-energy representation and why it complements the score.
%   Beat representation as model output -- contrast continuous beat confidence with discrete beat events, and explain why both are useful.
%   Onset visualization from MAPS -- explain how performance onsets are used as score-aligned markers.
%   Shared visualization context -- describe how the layers are combined on a common timeline and why this is necessary for comparison.
%   Design trade-off -- explain why some layers are kept vector-based while others are better represented as rasterized output.
%   End result -- describe the type of composite visualization the reader should expect at this stage, without going into low-level implementation details.


\section{Score and Data Alignment}
% Core narrative: 
%   Bridge matplotlib visualizations to the warped score's coordinate system.
% Key Topics:
%   Role of the warped score -- explain that the score SVG is the temporal reference that all other data must align to.
%   Timeline as the alignment backbone -- describe the score timeline as the structure that links score position to audio time.
%   Proportional scaling idea -- explain the simple rule-of-three relationship between score length and audio duration.
%   Why alignment matters -- show that alignment is what makes multi-layer visual comparison meaningful.
%   SVG as the common coordinate space -- explain that the score, onset markers, and embedded layers must share a compatible visual reference.
%   Raster layer placement -- mention that non-vector elements, such as spectrograms, need to be positioned using the same timing logic.
%   Reliability and limits -- note that the quality of alignment depends on the quality and density of the MAPS anchors.

    
In this section we'll go in depth on how the score and data visualization alignment functions at it's core. 



\section{Combining Visualizations}
% Core narrative: 
    % Show the two-stage export pipeline and multi-performance comparison.
% Key Topics:
%   Two-stage assembly process -- explain that visualizations are first created as separate layers and then merged into the aligned score.
%   Vector and raster combination -- describe how SVG layers and PNG-based layers coexist in the final composition.
%   Layer merging strategy -- explain that the final output is built by inserting the visual elements into the score structure rather than redrawing everything from scratch.
%   Multi-visualization comparison -- describe how multiple aligned outputs can be combined for side-by-side or stacked comparison.
%   Output as a publication-ready artifact -- emphasize that the final product is a reusable, high-resolution figure suitable for analysis and presentation.
%   Architectural takeaway -- close the section by showing that the system is designed for modular composition, not a single fixed visualization.
%   Bridge to the example walkthrough -- point forward to a concrete example script that demonstrates the full pipeline end to end.

\section{Conclusion}